% PDF visual theme for the statistical-analysis-course book.
% It mirrors the palette and visual language from assets/styles/book.css,
% while keeping the PDF close to the supplied reference edition.

\usepackage[most]{tcolorbox}
\usepackage{scrlayer-scrpage}
\usepackage{tikz}
\usepackage{caption}
\usepackage{colortbl}
\usepackage{array}
\usepackage{enumitem}
\usepackage{etoolbox}
\usepackage{ragged2e}

% -----------------------------------------------------------------------------
% Palette (same core colors as book.css)
% -----------------------------------------------------------------------------
\definecolor{BookInk}{HTML}{18262F}
\definecolor{BookInkSoft}{HTML}{42545F}
\definecolor{BookMuted}{HTML}{687984}
\definecolor{BookPaper}{HTML}{FFFFFF}
\definecolor{BookCanvas}{HTML}{F4F6F7}
\definecolor{BookLine}{HTML}{D9E0E3}
\definecolor{BookLineStrong}{HTML}{BCC9CE}
\definecolor{BookTeal}{HTML}{0F6972}
\definecolor{BookTealDark}{HTML}{094B52}
\definecolor{BookTealSoft}{HTML}{E7F1F1}
\definecolor{BookCoral}{HTML}{C7472F}
\definecolor{BookCoralSoft}{HTML}{F8ECE9}
\definecolor{BookGold}{HTML}{B97817}
\definecolor{BookGoldSoft}{HTML}{FBF3E4}
\definecolor{BookBlue}{HTML}{2A5EA8}
\definecolor{BookBiography}{HTML}{465B8C}
\definecolor{BookBiographySoft}{HTML}{F0F2F8}
% The HTML page uses a stronger sky blue. The reference PDF intentionally uses
% a paler fill for long derivations, so that equations remain comfortable to read.
\definecolor{BookBlueSoftPDF}{HTML}{EEF4FB}
\definecolor{BookFigureSoft}{HTML}{F9FBFB}
\definecolor{BookGallerySoft}{HTML}{FBFCFC}
\definecolor{BookTableStripe}{HTML}{F7F9FA}

% -----------------------------------------------------------------------------
% Base typography and paragraph rhythm
% -----------------------------------------------------------------------------
\AtBeginDocument{%
  \color{BookInk}%
  \KOMAoptions{parskip=false}%
  \setlength{\parindent}{1.2em}%
  \setlength{\parskip}{0.18\baselineskip}%
  \setlength{\emergencystretch}{3em}%
  \setcounter{secnumdepth}{3}%
  \setcounter{tocdepth}{0}%
}

\setlist[itemize]{leftmargin=1.55em,itemsep=0.14em,topsep=0.35em,parsep=0pt}
\setlist[enumerate]{leftmargin=1.7em,itemsep=0.14em,topsep=0.35em,parsep=0pt}

% -----------------------------------------------------------------------------
% Headings: sans serif, dark ink; chapter number in teal, like the reference PDF
% -----------------------------------------------------------------------------
\addtokomafont{disposition}{\sffamily\bfseries\color{BookInk}}
\setkomafont{chapter}{\sffamily\bfseries\fontsize{30}{33}\selectfont\color{BookInk}}
\setkomafont{section}{\sffamily\bfseries\fontsize{14.2}{17}\selectfont\color{BookInk}}
\setkomafont{subsection}{\sffamily\bfseries\fontsize{11.3}{14}\selectfont\color{BookInk}}
\setkomafont{subsubsection}{\sffamily\bfseries\fontsize{10.3}{13}\selectfont\color{BookInk}}

\renewcommand{\bookchaptercolor}{%
  \ifcase\value{part}BookTeal%
  \or BookTeal%
  \or BookGold%
  \or BookBlue%
  \or BookCoral%
  \else BookTeal\fi
}
\renewcommand{\BookChapterLabel}{Глава}
\renewcommand{\BookPartLabel}{РАЗДЕЛ КНИГИ}
\renewcommand{\BookPartContentsLabel}{Главы}
\renewcommand{\BookTocChapterLabel}{Глава}
\RedeclareSectionCommand[
  beforeskip=1.55\baselineskip,
  afterskip=0.45\baselineskip
]{section}
\RedeclareSectionCommand[
  beforeskip=1.15\baselineskip,
  afterskip=0.28\baselineskip
]{subsection}
\RedeclareSectionCommand[
  beforeskip=0.95\baselineskip,
  afterskip=0.2\baselineskip
]{subsubsection}

% -----------------------------------------------------------------------------
% Links and captions
% -----------------------------------------------------------------------------
\AtBeginDocument{%
  \renewcommand{\figurename}{Рис.}%
  \renewcommand{\tablename}{Таблица}%
  \renewcommand{\contentsname}{Содержание}%
}

% Pandoc can apply its own hyperlink setup after header includes.  Apply the
% book palette at the very end of \begin{document}, so external links such as
% "Слайды лекции" and "Задачи" are reliably teal in the PDF.
\AddToHook{begindocument/end}{%
  \hypersetup{
    colorlinks=true,
    linkcolor=BookTeal,
    citecolor=BookTeal,
    urlcolor=BookTeal,
    filecolor=BookTeal,
    linktoc=section
  }%
  \renewcommand{\contentsname}{Содержание}%
}

\captionsetup[figure]{
  font=small,
  labelfont=bf,
  textfont=it,
  labelsep=space,
  justification=centering,
  singlelinecheck=false,
  skip=7pt
}
\captionsetup[table]{
  font=small,
  labelfont=bf,
  labelsep=space,
  justification=raggedright,
  singlelinecheck=false
}

% Per-chapter bibliographies should remain with the chapter instead of leaving
% a nearly empty continuation page for a few long DOI and URL lines.
\AtBeginEnvironment{CSLReferences}{%
  \footnotesize
  \setlength{\emergencystretch}{2em}%
}

% -----------------------------------------------------------------------------
% Tables: light rules and alternating rows, matching the HTML and reference PDF
% -----------------------------------------------------------------------------
% IMPORTANT: xcolor/colortbl implements \rowcolors through global state. If that
% state is left active after a table, AMS matrix environments (pmatrix, bmatrix,
% etc.) can inherit the alternating row backgrounds. On a colored tcolorbox this
% appears as white/gray rectangles behind matrix entries (for example the n/k in
% a binomial coefficient). Reset the row-color machinery after every table.
\makeatletter
\newcommand{\BookResetRowColors}{%
  \global\@rowcolorsfalse
  \global\let\CT@row@color\relax
  \global\let\@rowcolors\@empty
}
\makeatother

\AtBeginEnvironment{longtable}{%
  \arrayrulecolor{BookLineStrong}%
  \rowcolors{2}{BookTableStripe}{white}%
  \setlength{\arrayrulewidth}{0.35pt}%
  \renewcommand{\arraystretch}{1.18}%
}
\AtEndEnvironment{longtable}{\BookResetRowColors}

\AtBeginEnvironment{tabular}{%
  \arrayrulecolor{BookLineStrong}%
  \rowcolors{2}{BookTableStripe}{white}%
  \renewcommand{\arraystretch}{1.18}%
}
\AtEndEnvironment{tabular}{\BookResetRowColors}

% -----------------------------------------------------------------------------
% Book boxes. These are PDF counterparts of the CSS classes.
% -----------------------------------------------------------------------------
\tcbset{
  bookbox/.style={
    enhanced,
    breakable,
    boxrule=0.35pt,
    arc=1.2mm,
    outer arc=1.2mm,
    left=3.7mm,
    right=3.7mm,
    top=2.8mm,
    bottom=2.8mm,
    before skip=9pt,
    after skip=9pt,
    fontupper=\normalfont,
    coltext=BookInk,
    shadow={0mm}{-0.7mm}{0.4mm}{black!5}
  }
}

\newtcolorbox{bookchapteropening}[1][]{
  bookbox,
  colback=BookCoralSoft,
  colframe=BookCoral!72,
  borderline west={2.6pt}{0pt}{BookCoral},
  title={#1},
  colbacktitle=BookCoralSoft,
  coltitle=BookCoral,
  fonttitle=\sffamily\bfseries\small,
  titlerule=0pt,
  top=2.0mm,
  toptitle=1.5mm,
  bottomtitle=0.4mm
}

\newtcolorbox{booknote}{
  bookbox,
  colback=BookGoldSoft,
  colframe=BookGold!70,
  borderline west={2.5pt}{0pt}{BookGold},
  boxrule=0.25pt,
  shadow={0mm}{0mm}{0mm}{white}
}

\newtcolorbox{bookderivation}{
  bookbox,
  colback=BookBlueSoftPDF,
  colframe=BookBlue!62,
  borderline west={2.6pt}{0pt}{BookBlue},
  title={Вывод},
  colbacktitle=BookBlueSoftPDF,
  coltitle=BookBlue,
  fonttitle=\sffamily\bfseries\small,
  titlerule=0pt,
  shadow={0mm}{0mm}{0mm}{white}
}

\newtcolorbox{bookchaptersummary}[1][Итоги главы]{
  bookbox,
  colback=BookTealSoft,
  colframe=BookTeal!62,
  borderline north={2.4pt}{0pt}{BookTeal},
  title={#1},
  colbacktitle=BookTealSoft,
  coltitle=BookTeal,
  fonttitle=\sffamily\bfseries\large,
  titlerule=0pt,
  shadow={0mm}{0mm}{0mm}{white}
}

\newtcolorbox{bookgalleryitem}{
  enhanced,
  breakable,
  colback=BookGallerySoft,
  colframe=BookLine,
  boxrule=0.35pt,
  arc=0.8mm,
  left=2.4mm,
  right=2.4mm,
  top=2.4mm,
  bottom=2.4mm,
  before skip=7pt,
  after skip=7pt,
  shadow={0mm}{0mm}{0mm}{white}
}

\newtcolorbox{bookapplet}[1][]{
  bookbox,
  colback=BookGallerySoft,
  colframe=BookLine,
  borderline north={2.8pt}{0pt}{BookTeal},
  title={#1},
  colbacktitle=BookGallerySoft,
  coltitle=BookInk,
  fonttitle=\sffamily\bfseries\large,
  titlerule=0pt
}

\newtcolorbox{bookbiography}{
  bookbox,
  colback=BookBiographySoft,
  colframe=BookBiography!60,
  borderline west={3.2pt}{0pt}{BookBiography},
  boxrule=0.25pt,
  breakable=false,
  shadow={0mm}{0mm}{0mm}{white}
}

% A reusable caption for figures which the Lua filter deliberately takes out of
% floating figure environments so that they can safely live inside tcolorboxes.
\newcommand{\BookInlineFigureCaption}[1]{%
  \par\vspace{0.45em}%
  \refstepcounter{figure}%
  {\centering\sffamily\fontsize{8.5}{10}\selectfont\color{BookMuted}%
    \figurename~\thefigure\enspace #1\par}%
}

% -----------------------------------------------------------------------------
% A reference-style title page. The text remains metadata-driven for title/author;
% subtitle/site can be changed in the two commands below.
% -----------------------------------------------------------------------------
\newcommand{\BookSubtitleText}{От вероятности к физическому результату}
\newcommand{\BookSiteText}{neutrinohit.github.io}
\newcommand{\BookBrandText}{NEUTRINOHIT}

\makeatletter
\newcommand{\BookMakeTitle}{%
  \begin{titlepage}
    \thispagestyle{empty}%
    \begin{tikzpicture}[remember picture,overlay]
      % teal spine line
      \fill[BookTeal] ([xshift=28.8mm]current page.north west)
        rectangle ([xshift=30.1mm]current page.south west);
      % coral/gold accent bar
      \fill[BookGold] ([xshift=30.1mm,yshift=-48mm]current page.north west)
        rectangle ([xshift=34.4mm,yshift=-98mm]current page.north west);
      % brand
      \node[anchor=north west,text=BookTeal,font=\sffamily\bfseries\small]
        at ([xshift=41mm,yshift=-21mm]current page.north west) {\BookBrandText};
      % hairline
      \draw[BookLine,line width=.45pt]
        ([xshift=37mm,yshift=-31.5mm]current page.north west) --
        ([xshift=-25mm,yshift=-31.5mm]current page.north east);
      % title
      \node[anchor=north west,text width=145mm,align=left,text=BookInk,
            font=\sffamily\bfseries\fontsize{31}{34}\selectfont]
        at ([xshift=41mm,yshift=-37mm]current page.north west) {\@title};
      % subtitle
      \node[anchor=north west,text width=145mm,align=left,text=BookMuted,
            font=\sffamily\fontsize{15}{18}\selectfont]
        at ([xshift=41mm,yshift=-70mm]current page.north west) {\BookSubtitleText};
      % author
      \node[anchor=north west,text width=145mm,align=left,text=BookInk,
            font=\rmfamily\fontsize{12}{15}\selectfont]
        at ([xshift=41mm,yshift=-96mm]current page.north west) {\@author};
      % site at bottom
      \node[anchor=south west,text=BookMuted,font=\sffamily\scriptsize]
        at ([xshift=37mm,yshift=25mm]current page.south west) {\BookSiteText};
    \end{tikzpicture}%
    \null
  \end{titlepage}%
  % The reference edition places the printed TOC immediately after the cover.
  % We render it here explicitly and keep Quarto's automatic PDF TOC disabled,
  % avoiding both a missing TOC and accidental duplication.
  \clearpage
  \tableofcontents
  \clearpage
}
\AtBeginDocument{\let\maketitle\BookMakeTitle}
\makeatother

% -----------------------------------------------------------------------------
% TOC and chapter-entry color accents
% -----------------------------------------------------------------------------
\RedeclareSectionCommand[tocbeforeskip=0.12\baselineskip]{chapter}
\RedeclareSectionCommand[tocbeforeskip=0.55\baselineskip]{part}

\AtBeginDocument{%
  \addtokomafont{partentry}{\sffamily\bfseries\color{BookTeal}}%
  \addtokomafont{partentrypagenumber}{\sffamily\bfseries\color{BookInk}}%
  \addtokomafont{chapterentry}{\sffamily\bfseries\color{BookTeal}}%
  \addtokomafont{chapterentrypagenumber}{\sffamily\bfseries\color{BookInk}}%
}
